\abstract
{Resumen}
{El libro analiza la implementación y evaluación de una red IoT basada en el estándar LoRaWAN, destacando su eficiencia para la transmisión de datos en aplicaciones de monitoreo ambiental. Además, describe la arquitectura de la red, incluyendo nodos, pasarelas y servidores, así como sus capas (física, MAC y aplicación). Igualmente, se detalla la selección e integración de sensores para medir la calidad del aire y el uso de \emph{software} libre para la gestión y visualización de datos. A través de un enfoque experimental, se diseñó, simuló e implementó la red en Bogotá, evaluando cobertura, interoperabilidad y rendimiento. Los resultados evidencian la viabilidad de LoRaWAN para soluciones escalables, de bajo consumo y aplicables a ciudades inteligentes.}
{Palabras clave}
{IoT, LoRaWAN, sensores, calidad del aire, redes inalámbricas, monitoreo.}

\noindent\rule{1.8cm}{0.4pt}

\vspace{3mm}

\begin{otherlanguage*}{english}
  \abstract
  {Abstract}
  {The book analyzes the implementation and evaluation of an IoT network based on the LoRaWAN standard, highlighting its efficiency for data transmission in environmental monitoring applications. It also describes the network architecture, including nodes, gateways, and servers, as well as its layers (physical, MAC, and application). Likewise, it details the selection and integration of sensors to measure air quality and the use of open-source \emph{software} for data management and visualization. Through an experimental approach, the network was designed, simulated, and implemented in Bogotá, evaluating coverage, interoperability, and performance. The results demonstrate the viability of LoRaWAN for scalable, low-power solutions applicable to smart cities.}
  {Keywords}
  {IoT, LoRaWAN, sensors, air quality, wireless networks, monitoring.}
\end{otherlanguage*}

\vspace{1cm}

\hfill\comocitar{Poveda Zafra,~J.\,N., Martín Santamaría,~L.\,E., Gaona García,~E.\,E., Tabares Amaya,~S.\,Y. y Bustos Pinzón,~A.\,F. (2026). Evaluación de una red IoT bajo el estándar LoRa. Universidad Distrital Francisco José de Caldas-Editorial UD.}{XXXXXXXXXX}
